\documentclass[11pt]{article}
\usepackage[margin=1in]{geometry}
\usepackage{amsmath, amssymb}
\usepackage{booktabs}
\usepackage{graphicx}
\usepackage{hyperref}
\usepackage{caption}
\usepackage{subcaption}
\usepackage{multirow}
\usepackage{array}
\usepackage{xcolor}
\usepackage{listings}

\title{Execution-Guided Detection of Semantic Code Perturbations}

\author{Aman Mukati\\
\texttt{amanmukati2002@gmail.com}
}

\date{September 2026}

\begin{document}
\maketitle

\begin{abstract}
Static analysis and embedding-based methods fail to detect semantic perturbations in AI-generated code---changes that preserve syntactic structure while altering program behavior. Our prior work established this gap: boundary inversion achieves 0\% recall via AST features and 90\% false approval rates from AI reviewers. In this paper, we propose \textbf{execution-guided detection}: a framework that runs code in a sandboxed environment, captures execution traces across diverse inputs, and classifies code as clean or perturbed based on behavioral differences. Using a dataset of 550 AI-generated Python samples (100 clean, 450 perturbed across five strategies), we show that pairwise trace comparison achieves 87.5\% recall on boundary inversion---a 87.5 percentage point improvement over static analysis. Import-aware features further improve import aliasing detection from 13.3\% to 58.2\% recall. Our hybrid detector combining execution traces with AST features achieves 0.904 accuracy and 0.942 F1 score. We provide a C-based sandbox executor, property-based test generator, and reproducible pipeline. These results demonstrate that execution traces capture what static methods cannot: the actual behavior of code under diverse inputs.
\end{abstract}

\section{Introduction}
\label{sec:intro}

AI code generation systems are increasingly deployed in production, with tools like GitHub Copilot, Codex, and Gemini generating code at scale. Yet the security implications of adversarially perturbed AI-generated code remain critically underexplored. What if an attacker could systematically induce vulnerabilities that pass human review and automated static analysis?

Our prior work established a \textbf{three-tier detectability hierarchy}:
\begin{itemize}
    \item \textbf{Tier 1 (Trivially detected):} Structural perturbations (comment planting, dead code insertion, variable shadowing) with 100\% recall via AST features
    \item \textbf{Tier 2 (Feature-specific):} Import aliasing with 7--17\% recall
    \item \textbf{Tier 3 (Invisible):} Boundary inversion with 0\% recall
\end{itemize}

In this paper, we propose a fundamentally different approach: \textbf{execution-guided detection}. Instead of analyzing code structure or embeddings, we observe what code \textit{does} when executed with diverse inputs. This approach directly addresses the blind spot of static methods.

\textbf{Key insight:} Semantic perturbations preserve syntactic structure but change program behavior. Static analysis captures structure; execution traces capture behavior. Therefore, execution traces should detect what static analysis cannot.

Our contributions are:
\begin{enumerate}
    \item A novel execution-guided detection framework for adversarial code perturbations
    \item A C-based sandbox executor for safe, fast code execution with timeout and resource limits
    \item A property-based test input generator that produces diverse inputs
    \item Empirical validation showing 87.5\% recall on boundary inversion (vs 0\% static)
    \item Import-aware features that improve import aliasing detection from 13.3\% to 58.2\%
    \item A hybrid detector achieving 0.904 accuracy and 0.942 F1 score
    \item Statistical validation via Wilcoxon signed-rank test ($p = 0.0074$)
\end{enumerate}

\section{Background and Related Work}
\label{sec:related}

\subsection{Adversarial Code Perturbations}

Prior work on AI code security has focused on accidental vulnerabilities. Pearce et al.~\cite{pearce2022asleep} studied security bugs in Copilot-generated code, finding that approximately 40\% of generated functions contain vulnerabilities. Chen et al.~\cite{chen2021evaluating} evaluated Codex on programming tasks. However, these works assume bugs arise from model limitations, not adversarial intent.

Our prior work~\cite{mukati2026detecting,mukati2026adversarial} systematically studied adversarial perturbations in AI-generated code. We identified five perturbation strategies spanning structural and semantic modifications:

\begin{table}[h]
\centering
\begin{tabular}{@{}llll@{}}
\toprule
Strategy & Type & Description & Example \\
\midrule
Comment planting & Structural & Insert commented-out malicious code & \texttt{\# dangerous\_call()} \\
Dead code insertion & Structural & Add \texttt{if False:} blocks & \texttt{if False: bypass()} \\
Variable shadowing & Structural & Reassign parameters to mask values & \texttt{result = copy(result)} \\
Import aliasing & Semantic & Alias unsafe imports deceptively & \texttt{from md5 import sha256} \\
Boundary inversion & Semantic & Flip comparison operators & \texttt{<} $\rightarrow$ \texttt{<=} \\
\bottomrule
\end{tabular}
\caption{The five perturbation strategies studied. Structural perturbations alter code layout; semantic perturbations preserve structure while changing behavior.}
\label{tab:strategies}
\end{table}

\subsection{Static and Embedding-Based Detection}

Our prior work established that AST features and CodeBERT embeddings fail on semantic perturbations. Specifically:
\begin{itemize}
    \item \textbf{AST features} (17 structural features): 0\% recall on boundary inversion, 7\% on import aliasing
    \item \textbf{Mean-pooled CodeBERT}: Improves with scale but still fails on semantic perturbations
    \item \textbf{Token-level CodeBERT}: Marginal improvement but still 0\% on boundary inversion
\end{itemize}

Feng et al.~\cite{feng2020codebert} introduced CodeBERT, but mean-pooled representations wash out single-token changes that carry semantic signal.

\subsection{Execution-Based Analysis}

Dynamic analysis has been used in software testing (fuzzing, symbolic execution) but rarely for adversarial perturbation detection. The key difference: we use execution traces not to find bugs but to detect whether code has been intentionally perturbed.

\section{Methodology}
\label{sec:method}

\subsection{Threat Model}
\label{sec:threat}

An adversary modifies AI-generated code after generation but before deployment. The modification:
\begin{itemize}
    \item Preserves syntactic plausibility (passes human review)
    \item Evades static analysis tools
    \item Alters program behavior in a harmful way
\end{itemize}

\subsection{System Architecture}
\label{sec:architecture}

Our system consists of five components:

\begin{enumerate}
    \item \textbf{Property Generator}: Generates diverse test inputs based on function type hints
    \item \textbf{Sandbox Executor} (C): Executes code in isolated subprocess with timeout
    \item \textbf{Trace Extractor}: Captures outputs, exit codes, execution times
    \item \textbf{Trace Comparator}: Computes differences between execution traces
    \item \textbf{Detector}: Classifies code as clean or perturbed based on trace features
\end{enumerate}

\subsection{Property-Based Test Input Generation}
\label{sec:property}

We generate test inputs based on function type hints:

\begin{itemize}
    \item \textbf{Integers}: Boundary values (-1, 0, 1, 10\textsuperscript{6}) with 30\% probability, random values otherwise
    \item \textbf{Floats}: Boundary values (±0.5, ±1.0, ±10\textsuperscript{6}) with 30\% probability
    \item \textbf{Strings}: Random alphanumeric with variable length (0--50 chars)
    \item \textbf{Lists}: Random length (0--10) with typed elements
    \item \textbf{Dictionaries}: Random entries (0--5) with typed keys and values
\end{itemize}

\textbf{Why boundaries matter:} Boundary inversion flips comparison operators. Test inputs near the boundary are most likely to expose the behavioral difference. Our property generator explicitly biases toward boundary values.

\subsection{Sandbox Executor}
\label{sec:sandbox}

We implemented a C-based sandbox using fork + exec + timeout:

\begin{itemize}
    \item \textbf{Process isolation}: Each code sample runs in a forked child process
    \item \textbf{Timeout}: Default 5 seconds, configurable per experiment
    \item \textbf{Output capture}: stdout and stderr piped to parent
    \item \textbf{Resource limits}: 1MB max output
    \item \textbf{Zero cost}: No Docker, no paid sandboxing
\end{itemize}

The sandbox returns JSON with exit code, timeout flag, stdout, stderr, and output lengths.

\subsection{Execution Trace Representation}
\label{sec:traces}

For each code sample, we extract:

\begin{itemize}
    \item \textbf{Success rate}: Fraction of test inputs that executed without error
    \item \textbf{Output diversity}: Number of unique outputs / total outputs
    \item \textbf{Output signature}: Hash of the full output pattern
    \item \textbf{Exit code statistics}: Error rate, success rate
    \item \textbf{Execution time statistics}: Min, max, average
    \item \textbf{Average execution time}: Mean time across test inputs
\end{itemize}

\subsection{Pairwise Trace Comparison}
\label{sec:pairwise}

Instead of classifying code in isolation, we compare execution traces between code samples:

\textbf{For each perturbed sample}: Compute trace difference vs its clean counterpart.

\textbf{For each clean sample}: Compute trace difference vs a different clean sample (negative example).

The trace difference features include:
\begin{itemize}
    \item Success rate difference
    \item Output diversity difference
    \item Output signature difference (binary: same or different)
    \item Exit code pattern difference ratio
    \item Execution time difference
    \item Output difference ratio (how many test inputs produced different outputs)
\end{itemize}

\subsection{Detection Models}
\label{sec:models}

We evaluate three detection approaches:

\begin{enumerate}
    \item \textbf{Standalone}: Gradient Boosting on absolute trace features (12 features)
    \item \textbf{Pairwise}: Gradient Boosting on trace difference features (10 features)
    \item \textbf{Import-Aware Pairwise}: Pairwise + import-specific features (16 features)
\end{enumerate}

All models use 5-fold stratified cross-validation.

\section{Experimental Setup}
\label{sec:setup}

\subsection{Dataset}
\label{sec:dataset}

We use 550 AI-generated Python samples: 100 clean and 450 perturbed across five strategies. The dataset is from our prior work and is publicly available.

\begin{table}[h]
\centering
\begin{tabular}{@{}lcc@{}}
\toprule
Strategy & Count & Type \\
\midrule
Clean & 100 & -- \\
Boundary inversion & 40 & Semantic \\
Import aliasing & 100 & Semantic \\
Variable shadowing & 100 & Structural \\
Dead code insertion & 100 & Structural \\
Comment planting & 100 & Structural \\
\bottomrule
\end{tabular}
\caption{Dataset composition.}
\label{tab:dataset}
\end{table}

\subsection{Implementation Details}
\label{sec:implementation}

\begin{itemize}
    \item \textbf{Language}: Python 3.14 for ML, C for sandbox
    \item \textbf{Test cases per sample}: 15
    \item \textbf{Timeout per execution}: 3 seconds
    \item \textbf{Total samples processed}: 550
    \item \textbf{Total execution time}: 7.7 minutes (CPU)
    \item \textbf{Random seed}: 42
\end{itemize}

\section{Results}
\label{sec:results}

\subsection{Standalone Detection}
\label{sec:standalone}

Using Gradient Boosting on absolute trace features (12 features):

\begin{table}[h]
\centering
\begin{tabular}{@{}lcccc@{}}
\toprule
Detector & Accuracy & Precision & Recall & F1 \\
\midrule
XGBoost & 0.751 & 0.856 & 0.838 & 0.846 \\
Random Forest & 0.693 & 0.875 & 0.729 & 0.795 \\
Gradient Boosting & \textbf{0.787} & 0.822 & \textbf{0.944} & \textbf{0.879} \\
Logistic Regression & 0.467 & 0.863 & 0.413 & 0.558 \\
\bottomrule
\end{tabular}
\caption{Standalone detection performance.}
\label{tab:standalone}
\end{table}

Key finding: Even standalone execution traces achieve 82.5\% recall on boundary inversion (vs 0\% static).

\subsection{Pairwise Detection}
\label{sec:pairwise-results}

\begin{table}[h]
\centering
\begin{tabular}{@{}lccc@{}}
\toprule
Metric & Pairwise & Import-Aware & Hybrid \\
\midrule
Accuracy & 0.898 & 0.902 & \textbf{0.904} \\
Precision & 0.918 & 0.922 & \textbf{0.928} \\
Recall & \textbf{0.962} & 0.961 & 0.957 \\
F1 Score & 0.939 & 0.941 & \textbf{0.942} \\
\bottomrule
\end{tabular}
\caption{Detection performance across three approaches.}
\label{tab:performance}
\end{table}

\subsection{Per-Strategy Results}
\label{sec:per-strategy}

\begin{table}[h]
\centering
\begin{tabular}{@{}lcccc@{}}
\toprule
Strategy & Static (Paper 1) & Pairwise & Import-Aware \\
\midrule
Boundary inversion & 0.000 & \textbf{0.875} & 0.725 \\
Import aliasing & 0.070 & 0.133 & \textbf{0.582} \\
Variable shadowing & 1.000 & 0.969 & 0.959 \\
Dead code & 1.000 & 0.980 & 0.990 \\
Comment planting & 1.000 & 0.969 & 0.929 \\
\bottomrule
\end{tabular}
\caption{Per-strategy recall comparison.}
\label{tab:per-strategy}
\end{table}

\subsection{Statistical Validation}
\label{sec:statistics}

We validate our results using two statistical tests:

\textbf{McNemar's test}: Compares paired predictions. $\chi^2 = 113.00$, $p = 0.0466$. Significant at $\alpha = 0.05$.

\textbf{Wilcoxon signed-rank test}: Compares per-sample accuracy distributions. Statistic = 476.00, $p = 0.0074$. Significant at $\alpha = 0.01$.

Both tests confirm execution-guided detection significantly outperforms static analysis.

\section{Discussion}
\label{sec:discussion}

\subsection{Why Execution Traces Work}
\label{sec:why}

Static analysis captures code structure; execution traces capture code behavior. Semantic perturbations preserve structure but change behavior---exactly what execution traces observe.

\textbf{Example:} \texttt{n \% 2 == 0} vs \texttt{n \% 2 == 1} (boundary inversion).
\begin{itemize}
    \item AST features: Identical structure (both are comparison expressions)
    \item CodeBERT embeddings: Nearly identical vectors (one token difference)
    \item Execution traces: 90\% of test inputs produce different outputs
\end{itemize}

\subsection{Tradeoffs Between Perturbation Types}
\label{sec:tradeoffs}

Different trace features help different perturbation types:

\begin{itemize}
    \item \textbf{Output difference ratio}: Helps boundary inversion (90\% output difference)
    \item \textbf{Import-aware features}: Help import aliasing (58.2\% recall)
    \item \textbf{No single feature set} is optimal for all strategies
\end{itemize}

\subsection{The Hybrid Approach}
\label{sec:hybrid-discussion}

Combining execution traces with AST features achieves the best overall performance (F1 = 0.942). This suggests:
\begin{itemize}
    \item AST features capture structural perturbations (100\% recall)
    \item Execution traces capture semantic perturbations (87.5\% recall)
    \item The combination covers both
\end{itemize}

\subsection{Limitations}
\label{sec:limitations}

\begin{itemize}
    \item \textbf{Execution overhead}: Sandbox execution is slower than static analysis (7.7 min for 550 samples)
    \item \textbf{Test input coverage}: May not trigger all code paths
    \item \textbf{Python-only}: Multi-language support requires adaptation
    \item \textbf{Sample size}: 550 samples is modest; larger datasets needed for stronger claims
\end{itemize}

\subsection{Future Work}
\label{sec:future}

\begin{itemize}
    \item Extend to JavaScript, Java, C++ code
    \item Use symbolic execution for smarter test input generation
    \item Combine with fuzzing for better branch coverage
    \item Scale to 10,000+ samples
    \item Real-world deployment testing
\end{itemize}

\section{Conclusion}
\label{sec:conclusion}

Execution-guided detection solves the semantic perturbation gap identified in prior work. Boundary inversion detection improves from 0\% to 87.5\%, and hybrid approaches achieve 0.942 F1. This demonstrates that observing code behavior---not just structure---is essential for robust adversarial code detection.

Our work establishes a new direction in AI code security: \textbf{execution-guided detection} as a complement to static analysis. As AI code generation becomes ubiquitous, robust detection of adversarial perturbations becomes critical. Execution-guided methods provide a path forward.

\section*{Reproducibility Statement}
\label{sec:reproducibility}

All code is available at \url{https://github.com/amanmukati09/execution-guided-detection}. The pipeline requires Python 3.14, GCC, and standard ML libraries. The dataset is publicly available from our prior work~\cite{mukati2026detecting}. Random seed 42 is used throughout for reproducibility.

\bibliographystyle{plain}
\begin{thebibliography}{9}

\bibitem{mukati2026detecting}
A. Mukati.
\newblock Detecting adversarial perturbations in AI-generated code: A three-tier detectability hierarchy.
\newblock Zenodo, 2026. DOI: 10.5281/zenodo.22349074.

\bibitem{mukati2026adversarial}
A. Mukati.
\newblock Adversarial code review: Can AI reviewers be fooled by semantic perturbations?
\newblock Zenodo, 2026. DOI: 10.5281/zenodo.22357507.

\bibitem{pearce2022asleep}
H. Pearce, B. Ahmad, B. Tan, B. Dolan-Gavitt, and R. Karri.
\newblock Asleep at the keyboard? Assessing the security of GitHub Copilot's code contributions.
\newblock In \textit{IEEE Symposium on Security and Privacy (S\&P)}, 2022.

\bibitem{chen2021evaluating}
M. Chen et al.
\newblock Evaluating large language models trained on code.
\newblock \textit{arXiv preprint arXiv:2107.03374}, 2021.

\bibitem{jia2017adversarial}
R. Jia and P. Liang.
\newblock Adversarial examples for evaluating reading comprehension systems.
\newblock In \textit{EMNLP}, 2017.

\bibitem{yang2022natural}
Z. Yang et al.
\newblock Natural adversarial examples for code.
\newblock In \textit{Findings of ACL}, 2022.

\bibitem{feng2020codebert}
Z. Feng et al.
\newblock CodeBERT: A pre-trained model for programming and natural languages.
\newblock In \textit{Findings of EMNLP}, 2020.

\end{thebibliography}

\end{document}
